\section{The joint information needed for a canonical profile}\label{sec:joint-source}

Source approximation and statistical evaluation answer different questions.
The preceding theorem controls the error in replacing the entire source.
We now specify the remaining joint input and give a finite accuracy
interface for it. The derivation follows the internally assessed
fixed-center construction in \cite{OCJoint,OCJointReview}.

\subsection{Fix the law before estimating it}

For each fixed finite $T\ge10$, use the compact window
\eqref{eq:full-window} with $h=1$, and set
\begin{equation}\label{eq:fixed-regime}
 \begin{gathered}
 X\sim T^{-1}\mathbf1_{[T,2T]}\dd X,\qquad
 \ell=\min(1,4\pi/\log T),\qquad Q=[-1,1]+\ii[-1,1],\\
 \tau=\tfrac14,\qquad\epsilon=\tau\ell,\qquad J=\ell^2.
 \end{gathered}
\end{equation}
Write $r_X,d_X,R_{2,X}$ for the unfolded first intensity, self term
and factorial pair density. Thus $R_{2,X}=r_X\otimes r_X-d_X$ and
$C_T=\mathbb E R_{2,X}-(\mathbb E r_X)\otimes(\mathbb E r_X)$.
All scales are positive and fixed at each $T$; this convention alone
asserts no limiting density as $T\to\infty$.

Use common source normalization $\mathsf M=4T+26$ and put
\[
 a_* =\log(2\mathsf M/c_0),\quad
 n=\max\{1,\lceil a_*/\log(27/26)\rceil\},\quad
 m=2n,\quad A_\epsilon=(2\pi\epsilon^2)^{-1}.
\]
Here $c_0$ is the anchor of Section~\ref{sec:canonical}.
The two Jensen disks bound the number of active labels by $m$ for
both the exact and an admissible approximating field. Consequently
\begin{equation}\label{eq:joint-population}
 0\le r_X\le R:=JmA_\epsilon,\qquad
 0\le d_X\le D:=J^2mA_\epsilon^2.
\end{equation}
These size bounds do not evaluate the covariance.

There is a useful exact frequency description. For
$a_j(X)=((x_j-X)/\ell,y_j/\ell)$ define
\[
 A_X(k)=\sum_jW_X(z_j)\ee^{-\ii k\cdot a_j(X)},\qquad
 B_X(k)=\sum_jW_X(z_j)^2\ee^{-\ii k\cdot a_j(X)}.
\]
With Fourier convention $\widehat f(k)=\int\ee^{-\ii k\cdot u}f(u)\dd^2u$,
Gaussian transformation gives
\begin{equation}\label{eq:joint-fourier}
 \widehat C_T(k,q)=\ee^{-\tau^2(|k|^2+|q|^2)/2}
 \bigl[\mathbb E A_X(k)A_X(q)-\mathbb E A_X(k)\mathbb E A_X(q)
                          -\mathbb E B_X(k+q)\bigr].
\end{equation}
The covariance is bilinear, without complex conjugation. Its absolute
value is at most $m^2$ by centered Cauchy--Schwarz. Restricting both
frequency vectors to disks of radius $K$ therefore gives inverse-transform
error at most
\begin{equation}\label{eq:joint-frequency-tail}
 \frac{m^2+m}{4\pi^2\tau^4}(2\ee^{-s}-\ee^{-2s}),
 \qquad s=\tau^2K^2/2.
\end{equation}
Indeed the mass of a two-dimensional Gaussian outside such a disk is
$\ee^{-s}$ times its total mass $2\pi/\tau^2$; the complement of the
product of two disks has relative mass $2\ee^{-s}-\ee^{-2s}$.

The mixed product retains the factor
\[
 \frac1T\int_T^{2T}W_X(z_j)W_X(z_k)
             \ee^{\ii(k_x+q_x)X/\ell}\dd X
\]
and the separate marks $k_y y_j+q_y y_k$. A complex-difference projection
does not specify this general four-real-frequency transform. Translating
two labels together preserves their mutual difference, but can change
the center overlap at an endpoint of $[T,2T]$. This is an information
test on finite configurations, not another canonical zero set. The
arithmetic projection actually available is retained in
Appendix~\ref{sec:arithmetic}.

The center edges also have a quantitative cost. Replacing the center
law by its normalized restriction to $[T+1,2T-1]$ changes $C_T$ by at
most $2(3R^2+D)/T$: the deleted probability is $2/T$, the raw product
costs $R^2$ times this probability, the product of means costs $2R^2$,
and the self mean costs $D$. Deleting labels outside $[T,2T]$ gives
the same bound, as a separate comparison, since those labels can
contribute only on the edge strips. Keeping the original law incurs
neither cost.

\subsection{Regularize the logarithm, keeping the nonharmonic part}

Translate the reflected disks to $D_\sigma=D(\sigma4\ii,r)$,
$r=25/6$, $\sigma\in\{+,-\}$, and let
$f_{X,\sigma}(\xi)=F_\sigma(X+\xi)/\mathsf M$. Set
\begin{align}
 \mathcal K_{\sigma,u}(\xi)
  &=\frac{J}{2\pi}\Delta_\xi
       [\theta_\sigma W_0G_\epsilon(\ell u-\xi)],\label{eq:joint-kernel}\\
 \mathcal D_{\sigma,u,v}(\xi)
  &=\frac{J^2}{2\pi}\Delta_\xi
       [\theta_\sigma W_0^2G_\epsilon(\ell u-\xi)
                                     G_\epsilon(\ell v-\xi)].\nonumber
\end{align}
They have compact support strictly inside their disks, are independent
of $X$, and contract $\log|f_{X,\sigma}|$ to $r_X$ and $d_X$.

\begin{proposition}[Scope of harmonic replacement]\label{prop:harmonic-joint}
If a reference field is harmonic on a neighborhood of each closed
disk $D_\sigma$, every contraction of that field with the corresponding
kernel in \eqref{eq:joint-kernel} vanishes. This holds pathwise for a
random harmonic reference, including all covariance cross terms.
\end{proposition}
\begin{proof}
Twice integrating by parts moves the compact test's Laplacian onto
the reference field. Its Laplacian is zero. Thus subtracting the
reference leaves the entire contracted statistic in the residual.
\end{proof}
This proposition concerns harmonic replacements on these full planar
supports. It does not exclude contour access, nonharmonic approximations
or estimates retaining the residual. In particular the real part of
a finite holomorphic prime polynomial has zero contraction; its leading
law cannot by itself determine this planar zero statistic.

For an integer $M\ge1$ define
\[
 P_M(t)=-\frac12\sum_{a=1}^M\frac{(1-t)^a}{a},\qquad
 Y_{X,\sigma}=1-|f_{X,\sigma}|^2\in[0,1],\qquad
 L_{X,\sigma}^{(M)}=P_M(|f_{X,\sigma}|^2).
\]
The sublevel estimate \eqref{eq:sublevel}, after squaring its threshold,
has the form
$|\{|f|^2<u\}|\le\Lambda u^{1/n}$, where
\[
 \Lambda=n\pi(26/3)^2\ee^{102a_*/n},\qquad \beta=n\Lambda/2.
\]
It extends to $u\ge1$ since $\Lambda$ exceeds the disk area.
Layer cake yields
$\int\ee^{-a|f|^2}\le\Lambda\Gamma(1+1/n)a^{-1/n}$.
The gamma factor is at most one by concavity of $x^{1/n}$ under
the unit exponential probability measure. Since $(1-t)^a\le\ee^{-at}$,
the nonnegative logarithm-series remainder satisfies
\begin{equation}\label{eq:polynomial-log-error}
 \int_{D_\sigma}|L_{X,\sigma}^{(M)}-\log|f_{X,\sigma}||\dd^2\xi
 \le\frac\Lambda2\sum_{a>M}a^{-1-1/n}
 \le\beta M^{-1/n}.
\end{equation}
Zeros cause no pointwise product estimate here; their contribution is
handled by integrability and the sublevel bound. The approximation is
nonharmonic: direct differentiation gives
$\Delta P_M(|f|^2)=2M|f'|^2(1-|f|^2)^{M-1}\ge0$.

Abbreviate $U=U_{1,\epsilon}$, $V=V_{1,\epsilon}$ from
\eqref{eq:UV-costs}. Contracting the two disk errors gives
\[
 \|r_X-r_X^{(M)}\|_\infty\le d_M:=\frac{JU\beta}{\pi}M^{-1/n},
 \qquad
 \|d_X-d_X^{(M)}\|_\infty\le e_M:=\frac{J^2V\beta}{\pi}M^{-1/n}.
\]
Each product of first intensities costs $2Rd_M+d_M^2$. Hence
\begin{equation}\label{eq:joint-bias}
 \|C_T-C_T^{(M)}\|_\infty\le4Rd_M+2d_M^2+e_M,
 \qquad
 \|\overline R_2-\overline R_2^{(M)}\|_\infty
       \le2Rd_M+d_M^2+e_M.
\end{equation}

\subsection{A finite mixed-moment accuracy theorem}

Define the actual same-center moments
\[
 \mu_a^\sigma(\xi)=\mathbb E Y_{X,\sigma}(\xi)^a,\qquad
 \mu_{ab}^{\sigma\upsilon}(\xi,\omega)
   =\mathbb E Y_{X,\sigma}(\xi)^aY_{X,\upsilon}(\omega)^b.
\]
Finite expansion gives the exact surrogate
\begin{align}
 C_T^{(M)}(u,v)
  ={}&\frac14\sum_{a,b=1}^M\frac1{ab}\sum_{\sigma,\upsilon}
     \iint\mathcal K_{\sigma,u}(\xi)\mathcal K_{\upsilon,v}(\omega)
       [\mu_{ab}^{\sigma\upsilon}-\mu_a^\sigma\mu_b^\upsilon]
                                      \dd^2\xi\dd^2\omega\nonumber\\
   &+\frac12\sum_{a=1}^M\frac1a\sum_\sigma
        \int\mathcal D_{\sigma,u,v}(\xi)\mu_a^\sigma(\xi)\dd^2\xi.
                                                        \label{eq:joint-moments}
\end{align}
The plus sign on the last line comes from subtracting the negative
polynomial logarithm in the self term; the kernel is signed. The
direct full-pair formula replaces the bracket by $\mu_{ab}^{\sigma\upsilon}$.

\begin{theorem}[Finite joint-source interface]\label{thm:joint-interface}
Put $H_M=\sum_{a=1}^M a^{-1}$, $\kappa=Jr^2U$ and $\nu=J^2r^2V$.
Suppose all first moments and mixed moments in
\eqref{eq:joint-moments} are supplied with uniform errors
$\delta_1,\delta_2$, and clip estimated first moments to $[0,1]$.
The additional connected-profile error is at most
\begin{equation}\label{eq:moment-error}
 \frac{\kappa^2H_M^2}{4}(\delta_2+2\delta_1)
                 +\frac{\nu H_M}{2}\delta_1.
\end{equation}
For the direct full pair the first parenthesis can be replaced by
$\delta_2$. For a complete-source approximation with common logarithm
error $E$, the connected source error is at most
\begin{equation}\label{eq:joint-source-error}
 \frac{J^2E}{\pi}(8nA_\epsilon U+V).
\end{equation}
Adding \eqref{eq:joint-source-error}, the first bound in
\eqref{eq:joint-bias}, and \eqref{eq:moment-error} gives a valid
total connected-profile error.
\end{theorem}
\begin{proof}
The sums of the two kernels' $L^1$ norms are bounded by $\kappa$ and
$\nu$: each disk has area $\pi r^2$ and the factors in
\eqref{eq:joint-kernel} are $J/(2\pi)$ and $J^2/(2\pi)$.
Products of clipped means change by at most $2\delta_1$.
Absolute integration of the finite sums proves \eqref{eq:moment-error}.
For \eqref{eq:joint-source-error}, the unfolded first-intensity error
is $JEU/\pi$ and both intensities are bounded by $2nJA_\epsilon$.
The raw product and the product of means each cost
$4nJ^2A_\epsilon EU/\pi$; the self term costs $J^2EV/\pi$.
\end{proof}

At fixed $T$, the complete arithmetic and radial tails can make $E$
arbitrarily small, and then $M$ can make \eqref{eq:joint-bias}
arbitrarily small. The remaining input is the same-center moment array,
or a stronger estimate for its contracted combination, at the tolerance
in \eqref{eq:moment-error}. Separate one-point bounds do not supply it.
Neither a practical complexity bound nor positivity of the finite-$M$
surrogate is asserted. This theorem specifies what an acquisition must
pay; it does not report an acquired canonical profile.
